% ch-a4-02.tex —— 第 4 篇（Emerton–Gee–Hellmann）续：§2–§3，底本 PDF 192–203 = 印刷页 179–190
\ysrc{179}
我们要感谢许多同事、合作者与朋友：我们曾与他们讨论本讲义所呈现的想法。特别要感谢 Johannes Anschütz、David Ben-Zvi、Roman Bezrukavnikov、George Boxer、Christophe Breuil、Frank Calegari、Ana Caraiani、Pierre Colmez、Gabriel Dospinescu、Andrea Dotto、Laurent Fargues、Tony Feng、Dennis Gaitsgory、David Geraghty、Ian Grojnowski、David Hansen、Michael Harris、David Helm、Christian Johansson、Arthur-César Le Bras、Bao Le Hung、Brandon Levin、Jacob Lurie、Lucas Mann、David Nadler、James Newton、Lue Pan、Vytautas Paškūnas、Vincent Pilloni、Alice Pozzi、Juan Esteban Rodriguez Camargo、Joaquin Rodrigues Jacinto、David Savitt、Peter Scholze、Benjamin Schraen、Jack Sempliner、Sug Woo Shin、James Timmins、Pol van Hoften、Akshay Venkatesh、Carl Wang-Erickson 与 Xinwen Zhu。

我们十分感谢审稿人提出的许多有益意见与更正。还要感谢 John Bergdall、Pierre Colmez、Fred Diamond、Andrea Dotto、Elmar Große-Klönne、David Hansen、Claudius Heyer、Heejong Lee、Vytautas Paškūnas 与 David Savitt 对本讲义早先版本的意见。

\section{记号}

\subsection{域}

固定 $\Q$ 的一个代数闭包 $\overline{\Q}$，并对每个素数 $p$ 固定 $\Q_p$ 的一个代数闭包 $\overline{\Q}_p$。若 $F/\Q_p$ 是有限扩张，记 $\mathrm{Gal}_F$ 为绝对伽罗瓦群 $\mathrm{Gal}(\overline{\Q}_p/F)$；记 $I_F$ 为 $\mathrm{Gal}_F$ 的惯性子群，$W_F$ 为 Weil 群，$k_F$ 为 $F$ 的整数环 $\OO_F$ 的剩余域。我们按如下方式正规化局部类域论：单值化元对应几何 Frobenius 元素。

若 $M$ 是数域，记 $\mathrm{Gal}_M:=\mathrm{Gal}(\overline{\Q}/M)$ 为 $M$ 的绝对伽罗瓦群。若 $S$ 是 $M$ 的一集位，则记 $M(S)/M$ 为在 $S$ 之外非分歧的最大扩域，并记 $\mathrm{Gal}_{M,S}:=\mathrm{Gal}(M(S)/M)$。对每个 $p$ 固定嵌入 $\overline{\Q}\hookrightarrow\overline{\Q}_p$，并固定嵌入 $\overline{\Q}\hookrightarrow\C$；于是若 $v$ 是 $M$ 的一个位，就有同态 $\mathrm{Gal}_{M_v}\to\mathrm{Gal}_M$。

记 $\OO$ 为 $\Q_p$ 的某个固定有限扩张 $L$ 的整数环，$k$ 为 $\OO$ 的剩余域，$\varpi$ 为 $\OO$ 的一个单值化元。（这些将充当各种表示的系数。）

\subsection{$p$-进 Hodge 理论}

若 $\rho$ 是 $\mathrm{Gal}_F$ 在 $L$-向量空间 $W$ 上的 de Rham 表示，则记 $\mathrm{WD}(\rho)$ 为对应的 $W_F$ 的 Weil--Deligne 表示（该构造的细节见 §5.2）；又若 $\sigma:F\hookrightarrow L$ 是域的连续嵌入，则记 $\mathrm{HT}_\sigma(\rho)$ 为 $\rho$ 关于 $\sigma$ 的 Hodge--Tate 权之重集（multiset），按定义其中 $i$ 的重数为 $\dim_L\bigl(W\otimes_{\sigma,F}\overline{F}(i)\bigr)^{\mathrm{Gal}_F}$；例如，若 $\varepsilon$ 表示 $p$-进分圆特征，则 $\mathrm{HT}_\sigma(\varepsilon)=\{-1\}$。

设 $L$ 包含一切连续嵌入 $F\hookrightarrow L$ 的像。\tn{底本此处作「$L$ 包含一切连续嵌入 $F\hookrightarrow L$ 的像」，按字面这句话是平凡成立的；据上下文（须让 $L$ 足够大以容纳一切嵌入的像）疑为「$L$ 包含一切连续嵌入 $F\hookrightarrow\overline{\Q}_p$ 的像」之误，译文照原。}所谓 $d$ 元带标号 Hodge--Tate 权（labeled Hodge--Tate weights）$\underline{\lambda}$，指的是整数组 $\{\lambda_{\sigma,i}\}_{\sigma:F\hookrightarrow L,\,1\le i\le d}$，其中对一切 $\sigma$ 与一切 $1\le i\le d-1$ 都有 $\lambda_{\sigma,i}\ge\lambda_{\sigma,i+1}$。我们也把 $\underline{\lambda}$ 称为一个 Hodge 型（Hodge type）。若对一切 $\sigma$ 与 $1\le i\le d-1$ 都有 $\lambda_{\sigma,i}>\lambda_{\sigma,i+1}$，则称 $\underline{\lambda}$ 正则。若对每个 $\sigma:F\hookrightarrow L$ 都有 $\mathrm{HT}_\sigma(\rho)=\{\lambda_{\sigma,i}\}_{1\le i\le d}$，则称 de Rham 表示 $\rho$ 具有 Hodge 型 $\underline{\lambda}$（或带标号 Hodge--Tate 权 $\underline{\lambda}$）。

\ysrc{180}
所谓惯性型（inertial type）$\tau$，指的是表示 $\tau:I_F\to\mathrm{GL}_d(L)$，它能延拓为 $W_F$ 的、核为开子群的表示（特别地，$\tau$ 的像有限）。若 $\mathrm{WD}(\rho)|_{I_F}\cong\tau$，则称 de Rham 表示 $\rho$ 具有惯性型 $\tau$。

\subsection{Hodge--Tate 权与最高权}

在 $p$-进局部 Langlands 对应之下，对应式伽罗瓦一侧的 Hodge--Tate 权将与表示论一侧的最高权相关联；备一套适应对应式两侧的记号是有用的。按惯例，Hodge--Tate 权与最高权相差一个「$\rho$-平移」（$\rho$-shift）。

\begin{defn}{2.3.1}\label{defn:a4-2-3-1}
若 $\underline{\lambda}$ 是正则 Hodge 型，则记
\[
\xi_{\sigma,i}=i-1-\lambda_{\sigma,d+1-i},
\]
于是 $\xi_{\sigma,1}\ge\cdots\ge\xi_{\sigma,d}$。
\end{defn}

\begin{defn}{2.3.2}\label{defn:a4-2-3-2}
若 $\underline{\lambda}$ 是正则 Hodge 型，则把每个 $\xi_\sigma:=(\xi_{\sigma,1},\dots,\xi_{\sigma,d})$ 视为代数群 $\mathrm{GL}_d$ 的支配权（dominant weight）（相对于上三角 Borel 子群）。令 $W_{\underline{\lambda}}$ 为 $\mathrm{GL}_d(\OO_F)$ 的相应 $\OO$-表示，其定义如下：对每个 $\sigma:F\hookrightarrow L$，记 $W_{\xi_\sigma}$ 为 $K$ 的最高权为 $\xi_\sigma$ 的代数 $\OO_F$-表示（更确切地说，对偶 Weyl 模（dual Weyl module））。然后定义
\[
W_{\underline{\lambda}}:=\bigotimes_\sigma W_{\xi_\sigma}\otimes_{\OO_F,\sigma}\OO,
\]
并记 $L(\xi)=W_{\underline{\lambda}}[1/p]$，它是 $\mathrm{GL}_d(F)$（或代数群 $\mathrm{Res}_{F/\Q_p}\mathrm{GL}_d$）的表示。
\end{defn}

注意：尽管记号如此，$W_{\underline{\lambda}}$ 的最高权对应的是 $\xi_\sigma$，而非 $\lambda_\sigma$。

\begin{rem}{2.3.3}\label{rem:a4-2-3-3}
上述记号未必最优，但足敷所用。「$W$」代表（对偶）Weyl；强调对 $\underline{\lambda}$（而非对 $\xi$）的依赖，便于与 $p$-进 Hodge 理论的其他考量相比较。记号 $L(\xi)$ 是照 $\mathcal{O}$ 范畴（category $\mathcal{O}$）理论中的传统选定的。事实上，对（不必支配、不必整的）任一权 $\xi$（属于 $\mathrm{Res}_{F/\Q_p}\mathrm{GL}_d$ 的李代数），我们用 $M(\xi)$ 记最高权为 $\xi$ 的 Verma 模，用 $L(\xi)$ 记 $M(\xi)$ 的唯一单商。当 $\xi$ 支配且整时，$L(\xi)$ 恰是通常的最高权表示 $W_{\underline{\lambda}}[1/p]$，这就解释了此处记号的选取。
\end{rem}

\section{作为范畴 $p$-进局部 Langlands 之动机的 Taylor--Wiles 拼接}

本节回顾 Taylor--Wiles 拼接方法与 $\mathrm{GL}_2(\Q_p)$ 的 $p$-进局部 Langlands 对应的若干历史，并尽可能强调二者之间的联系。我们不打算给出详尽综述，而是一律透过范畴 Langlands 纲领的视角来看。读者若想了解 Taylor--Wiles 方法自创立以来发展的更主流叙述，可参阅 Calegari 的综述 [Cal23]。

\subsection{拼接的极简导引}

拼接最初由 Taylor 与 Wiles 引入，作为证明模性提升定理的技术 [TW95]。此后 Kisin 进一步发展了它，说明它提供了把 $p$-进伽罗瓦表示理论的局部面与整体面相联系的机制 [Kis09c, Kis07a, Kis09a, Kis10b, GK14]。

\ysrc{181}
本节通篇假设 $p>2$。我们先置身于如下简单却富有启发的背景中。固定连续、绝对不可约表示 $\overline{r}:\mathrm{Gal}_{\Q,\{p,\infty\}}\to\mathrm{GL}_2(k)$，其行列式为 $\varepsilon^{-1}$，并设 $\overline{\rho}:=\overline{r}|_{\mathrm{Gal}_{\Q_p}}$ 与 $\overline{r}|_{\mathrm{Gal}_{\Q(\zeta_p)}}$ 也都绝对不可约。我们假设 $\overline{r}$ 是模的（等价地，由 Serre 猜想 [KW09]，假设 $\overline{r}$ 是奇的）。

对任何包含 $p$ 与 $\infty$ 的有限素数集 $S$，可以考虑 $\overline{r}$ 在 $\OO$ 上的整体形变空间 $X_S(\overline{r})$，它参数化 $\overline{r}$ 到完备诺特局部 $\OO$-代数上 $\mathrm{Gal}_{\Q,S}$ 的表示的形变，且行列式固定为 $\varepsilon^{-1}$。也可以考虑局部形变空间 $X(\overline{\rho})$，它参数化 $\overline{\rho}$ 到完备诺特局部 $\OO$-代数上 $\mathrm{Gal}_{\Q_p}$ 的表示的形变，行列式同样固定为 $\varepsilon^{-1}$。二者分别是某个完备诺特局部 $\OO$-代数的 $\Spf$，这两个代数分别记作 $R_S(\overline{r})$ 与 $R(\overline{\rho})$。伽罗瓦表示的限制诱导态射
\begin{equation}
\tag{3.1.1}
X_S(\overline{r})\to X(\overline{\rho}),
\end{equation}
它是有限态射（见 [EP20] 与 [AC14]）。

$X_S(\overline{r})$（相应地 $X(\overline{\rho})$）闭点处的约化切空间等于整体伽罗瓦上同调群 $H^1\bigl(G_{\Q,S},\mathrm{Ad}^0(\overline{r})\bigr)$（相应地等于局部伽罗瓦上同调群 $H^1\bigl(G_{\Q_p},\mathrm{Ad}^0(\overline{\rho})\bigr)$）；形变空间之间的限制映射诱导约化切空间上的限制映射
\begin{equation}
\tag{3.1.2}
H^1\bigl(G_{\Q,S},\mathrm{Ad}^0(\overline{r})\bigr)\to H^1\bigl(G_{\Q_p},\mathrm{Ad}^0(\overline{\rho})\bigr),
\end{equation}
它由上同调类通常的拉回给出。

为简化情形，我们假设当 $S=\{p,\infty\}$ 时 (3.1.1) 其实是闭浸入；等价地，假设 Selmer 群（Selmer group）$\mathrm{Sel}\bigl(G_{\Q,\{p,\infty\}},\mathrm{Ad}^0(\overline{r})\bigr)$ 消失（按定义，该 Selmer 群是 (3.1.2) 的核；即它由要求上同调类在 $p$ 处局部消失来定义）。在此假设下（以及通篇假设 $\overline{r}|_{\mathrm{Gal}_{\Q(\zeta_p)}}$ 绝对不可约），Greenberg--Wiles 公式（[DDT97, Thm. 2.18]）表明
\[
\dim_k H^1\bigl(G_{\Q,\{p,\infty\}},\mathrm{Ad}^0(\overline{r})(1)\bigr)=\dim_k \mathrm{Sel}\bigl(G_{\Q,\{p,\infty\}},\mathrm{Ad}^0(\overline{r})\bigr)+1=1.
\]
（此式左端出现的 $H^1$，在我们的语境中就是 Greenberg--Wiles 公式里的「对偶 Selmer 群」（dual Selmer group）。）Taylor--Wiles 方法于是要求我们考虑满足 $q\equiv1\bmod p$ 的素数 $q$，使得限制态射
\[
H^1\bigl(G_{\Q,\{p,\infty\}},\mathrm{Ad}^0(\overline{r})(1)\bigr)\to H^1\bigl(\mathrm{Gal}_{\Q_q},\mathrm{Ad}^0(\overline{r})(1)\bigr)
\]
非零（等价地，为单射），且 $\overline{r}(\mathrm{Frob}_q)$ 的特征值互不相同；这就是所谓的 Taylor--Wiles 素数（Taylor--Wiles prime）。再用一次 Greenberg--Wiles 公式便知，对这种 Taylor--Wiles 素数 $q$，
\[
\mathrm{Sel}\bigl(\mathrm{Gal}_{\Q,\{p,q,\infty\}},\mathrm{Ad}^0(\overline{r})\bigr)=0
\]
（Selmer 条件同样由要求上同调类在 $p$ 处局部消失来定义），从而形变空间的限制态射
\begin{equation}
\tag{3.1.3}
X_{\{p,q,\infty\}}(\overline{r})\to X(\overline{\rho})
\end{equation}
仍是闭浸入。

\ysrc{182}
记 $\overline{r}_q:=\overline{r}|_{\mathrm{Gal}_{\Q_q}}$。在假设 $q\equiv1\bmod p$ 且 $\overline{r}(\mathrm{Frob}_q)$ 特征值互异之下，$\overline{r}_q$ 的形变问题可表（representable）（尽管 $\overline{r}_q$ 的自同态多于标量）。此外，形变环 $R(\overline{r}_q)$ 的形状非常简单：$R(\overline{r}_q)$ 形如 $\OO[[x,y]]/((1+x)^{p^n}-1)$，其中 $p^n$ 是整除 $q-1$ 的 $p$ 的最大幂。（更确切地说，可以验证 $\overline{r}_q$ 的一切形变都是特征之和，其中每个特征分别形变 $\overline{r}_q$ 为其直和分解中的那两个非分歧特征之一。由于我们的形变行列式固定，其中一个特征决定另一个；于是，一次性固定 $\overline{r}_q$ 的两个直和项之一后即可看出，$R(\overline{r}_q)$ 与这个固定非分歧特征的形变环重合。变量 $x,y$ 则分别对应于驯顺惯性群（tame inertia group）的一个生成元与一个 Frobenius 元素。）

现在考虑形变空间的限制态射
\[
X_{\{p,q,\infty\}}(\overline{r})\to X(\overline{r}_q)
\]
（当然 $X(\overline{r}_q):=\Spf R(\overline{r}_q)$）。由上一段的讨论，可以把它改写成态射
\begin{equation}
\tag{3.1.4}
X_{\{p,q,\infty\}}(\overline{r})\to\Spf \OO[[x]]/((1+x)^{p^n}-1),
\end{equation}
它在轨迹 $x=0$ 上的纤维等于 $X_{\{p,\infty\}}(\overline{r})$；这样便可把 $X_{\{p,q,\infty\}}(\overline{r})$ 看作 $X_{\{p,\infty\}}(\overline{r})$ 的某种加厚（thickening）\footnote{因为若模去 $\OO$ 的单值化元，则 $\Spec k[[x]]/((1+x)^{p^n}-1)=\Spec k[x]/x^{p^n}$ 是 $\Spec k$ 的幂零加厚。（底本此处作 $\Spec k[[x]]/(1+x)^{p^n}-1)$，括号不配对，译文按显然的意思补全。）}。（若态射 (3.1.4) 平坦，这种解读方式就是最优的。这并不总成立；但 Taylor--Wiles 方法给出的结论是：$X_{\{p,q,\infty\}}(\overline{r})$ 承载一个在 $\OO[[x]]/((1+x)^{p^n}-1)$ 上平坦的忠实模，即下文引入的完备同调（completed homology）群 $\widetilde{H}_1(q)$，故情形总归相当接近最优。）

事实是，$X_{\{p,\infty\}}(\overline{r})$ 在 $\OO$ 上的相对维数为 2，而 $X(\overline{\rho})$ 在 $\OO$ 上的相对维数为 3。于是闭浸入 (3.1.1) 的余维为 1；闭浸入 (3.1.3) 则把 $X_{\{p,q,\infty\}}(\overline{r})$ 实现为 $X_{\{p,\infty\}}(\overline{r})$ 沿其在 $X(\overline{\rho})$ 中的横截方向的加厚（若参照 (3.1.4)，是沿「$x$ 方向」的加厚）。

关键的一点是：对 $p$ 的任一给定幂 $p^n$，都能找到大量 Taylor--Wiles 素数 $q$（归根结底靠 Čebotarev 论证），使 $q\equiv1\bmod p^n$。于是可以在 $X(\overline{\rho})$ 内把 $X_{\{p,q,\infty\}}(\overline{r})$ 任意加厚；拼接方法核心处那个颇令人意外的想法，就是试图把 $X(\overline{\rho})$ 实现为这许多加厚 $X_{\{p,q,\infty\}}(\overline{r})$ 的某种极限。

为使之精确，先把每个形变环 $R(\overline{r}_q)$ 视为 $S_\infty:=\OO[[x]]$-代数（这使上文所说的「$x$ 方向」随 $q$ 变化而有一致的含义），然后在 $S_\infty$ 上把各个形式概形 $X_{\{p,q,\infty\}}(\overline{r})$ 随 $q$ 变化粘起来，且粘法要与其到 $X(\overline{\rho})$ 的嵌入相容；等价地，把各个 $(S_\infty,R(\overline{\rho}))$-双代数 $R_{p,q,\infty}(\overline{r})$ 粘起来。作为拼接过程的一部分，我们同时把各个（完备）同调群也粘起来。

回忆我们把模曲线的完备（第一）同调定义为
\begin{equation}
\tag{3.1.5}
\widetilde{H}_1:=\varprojlim_r H_1\bigl(Y(p^r),\OO\bigr)\cong\varprojlim_r H^1_c\bigl(Y(p^r),\OO\bigr).
\end{equation}

\ysrc{183}
这里 $Y(p^r)$ 是水平 $p^r$ 的模曲线（更确切地说，是 $\mathrm{PGL}_2$-对称空间的某个同余商，即 $\mathrm{GL}_2$ 的通常模曲线的若干连通分支之并）。

$\widetilde{H}_1$ 上作用有由素于 $p$ 的 Hecke 算子生成的 Hecke 代数 $\mathbf T$，也作用有 $\mathrm{PGL}_2(\Q_p)$ 与 $\mathrm{Gal}_{\Q,\{p,\infty\}}$。这些作用彼此交换。$\mathbf T$ 有一个对应模伽罗瓦表示 $\overline{r}$ 的极大理想 $\mathfrak{m}$。把 $\widetilde{H}_1$ 在 $\mathfrak{m}$ 处局部化，得到 $\mathbf T_{\mathfrak{m}}$-模 $\widetilde{H}_{1,\mathfrak{m}}$。类似地，对上文每个 $q$ 有完备同调群
\[
\widetilde{H}_1(q):=\varprojlim_r H_1\bigl(Y(p^r q),\OO\bigr)\cong\varprojlim_r H^1_c\bigl(Y(p^r q),\OO\bigr),
\]
其上作用有 $\mathrm{PGL}_2(\Q_p)$、$\mathrm{Gal}_{\Q,\{p,q,\infty\}}$、素于 $pq$ 的 Hecke 算子，以及 $q$ 处的某些 Hecke 算子。我们可以（也确实）把它们在 $\mathfrak{m}$ 处局部化（此时 $\mathfrak{m}$ 已扩张以纳入 $q$ 处的一个 Hecke 算子；确切的扩张方式对应于上文给出 $R(\overline{r}_q)$ 显式描述时对 $\overline{r}_q$ 的两个直和项所作的选取），得到 $S_\infty$-模与 $R(\overline{\rho})$-模 $\widetilde{H}_{1,\mathfrak{m}}(q)$。事实上，$S_\infty$ 在 $\widetilde{H}_{1,\mathfrak{m}}(q)$ 上的作用按定义经由我们所选态射 $S_\infty\to R(\overline{r}_q)$ 与自然态射 $R(\overline{r}_q)\to R_{p,q,\infty}(\overline{r})$ 的复合；于是利用 $S_\infty$ 的形式光滑性（formal smoothness）与闭浸入 (3.1.3)，可以构造态射 $S_\infty\to R(\overline{\rho})$，它与两个环在 $\widetilde{H}_{1,\mathfrak{m}}(q)$ 上的作用都相容。

仿照 [Sch18, §9]，作超积（ultraproduct）构造，得到带 $\mathrm{PGL}_2(\Q_p)$ 作用的 $R(\overline{\rho})$-模 $M_\infty$，具有性质
\[
M_\infty\otimes_{R(\overline{\rho})}R_{p,\infty}(\overline{r})=\widetilde{H}_{1,\mathfrak{m}}.
\]
按构造（以及素数 $q$ 处的局部--整体相容性），$M_\infty$ 是有限投射 $S_\infty[[K]]$-模，其中 $K=\mathrm{PGL}_2(\Z_p)$。

一般并不能假设 Selmer 群 $\mathrm{Sel}\bigl(G_{\Q,\{p,\infty\}},\mathrm{Ad}^0(\overline{r})\bigr)$ 消失（等价地，不能假设 (3.1.1) 是闭浸入），此时 Taylor--Wiles 方法如下。令 $g:=\dim_k \mathrm{Sel}\bigl(G_{\Q,\{p,\infty\}},\mathrm{Ad}^0(\overline{r})\bigr)+1$。上文提到的 Čebotarev 论证表明：对每个 $n$ 都存在一个 Taylor--Wiles 素数集 $\mathcal{Q}_n$，使得
\begin{equation}
\tag{3.1.6}
\dim_k \mathrm{Sel}\bigl(G_{\Q,\{p,\infty\}\cup\mathcal{Q}_n},\mathrm{Ad}^0(\overline{r})\bigr)=\dim_k \mathrm{Sel}\bigl(G_{\Q,\{p,\infty\}},\mathrm{Ad}^0(\overline{r})\bigr),
\end{equation}
且对每个 $q\in\mathcal{Q}_n$ 有 $q\equiv1\bmod p^n$。记 $S_\infty:=\OO[[x_1,\dots,x_g]]$，$R_\infty:=R(\overline{\rho})[[y_1,\dots,y_{g-1}]]$，其中 $x_i,y_i$ 为形式变量。由 (3.1.6)，可以把 (3.1.1) 与 (3.1.3) 提升为（非典范的）闭浸入 $X_{\{p,\infty\}}(\overline{r})\to\Spf R_\infty$ 与 $X_{\{p,\infty\}\cup\mathcal{Q}_n}(\overline{r})\to\Spf R_\infty$；上文在 $g=1$ 情形所解释的构造经直接推广，给出带 $R_\infty$ 作用的有限投射 $S_\infty[[K]]$-模 $M_\infty$，并且同样可以构造与各环在 $M_\infty$ 上的作用都相容的态射 $S_\infty\to R_\infty$。按构造，仍有
\begin{equation}
\tag{3.1.7}
M_\infty\otimes_{R_\infty}R_{p,\infty}(\overline{r})=\widetilde{H}_{1,\mathfrak{m}}.
\end{equation}

拼接模 $M_\infty$ 可以看作从完备（上）同调中提取 $p$ 处局部信息的一种办法。其构造依赖许多选择，但 [CEG+16] 首次采取的视角是：它在显然的意义下应是纯局部的；而对更一般的志村簇作同样的构造，应能给出超越 $\mathrm{GL}_2(\Q_p)$ 情形的 $p$-进局部 Langlands 对应的候选。这一预期尚未获验证，但它在发现本讲义所呈现的猜想时对我们十分重要，而模 $M_\infty$ 也出现在我们猜想的陈述中（特别见注 6.1.31）。

\ysrc{184}
\subsubsection*{3.1.8.\ 有限层上的拼接}

若 $V$ 是任一有限生成 $\OO$-模，带 $K$ 的连续作用，则置
\begin{equation}
\tag{3.1.9}
M_\infty(V):=M_\infty\otimes_{\OO[[K]]}V.
\end{equation}
于是 $V\mapsto M_\infty(V)$ 是正合函子，且 (3.1.7) 蕴含
\[
M_\infty(V)\otimes_{R_\infty}R_{p,\infty}(\overline{r})=H^1_c(V)_{\mathfrak{m}},
\]
其中 $H^1_c(V)_{\mathfrak{m}}$ 记系数取在 $V$ 所决定的局部系统中的紧支上同调群（在 $\mathfrak{m}$ 处局部化）。（由于 $\mathfrak{m}$ 非 Eisenstein，自然映射 $H^1_c(V)_{\mathfrak{m}}\to H^1(V)_{\mathfrak{m}}$ 是同构；但我们沿用 $H^1_c$ 记号，以强调所处理的是同调。）

\begin{rem}{3.1.10}\label{rem:a4-3-1-10}
如下文所示，刚才用 $M_\infty$ 给出的 $M_\infty(V)$ 之定义在历史上是倒过来的。事实上，最初的 Taylor--Wiles 方法固定一个 $V$，并构造相应的拼接环——原始文献中记作 $R_\infty$，我们或许应记作 $R_\infty(V)$——回过头看，它可以解释为上文引入的环 $R_\infty$ 中在 $M_\infty(V)$ 上忠实作用的那个商。在 Taylor 与 Wiles 最初考虑的情形中，拼接模 $M_\infty(V)$ 在 $R_\infty$ 上是循环的，从而是 $R_\infty(V)$ 上的自由模，故只以隐含方式出现（体现在 mod $p$ 重数一（multiplicity one）结果的运用中：这些结果表明，由模曲线上同调产生的某些 Hecke 模在其上作用的 Hecke 代数上是自由的）。拼接模 $M_\infty(V)$ 本身由 Diamond [Dia97] 与 Fujiwara [Fuj06] 各自独立引入；此后 Kisin 的工作 [Kis09a]（经 [EGS15] 解读）把拼接重述为以正合函子 $V\mapsto M_\infty(V)$ 为出发点。该函子由单个「大」拼接模 $M_\infty$ 决定，这一点是 [CEG+16] 观察到的。（不过要注意，公式 (3.1.9) 并未逐字出现于 [CEG+16]；那里出现的是本质等价（但写法更繁）的、涉及 Hom 与对偶的公式。）
\end{rem}

读者不妨把拼接函子的公式 (3.1.9) 与猜想 6.1.15 (2) 对照，后者提示助记式「$M_\infty=M_\infty(\OO[[K]])$」。从考虑函子 $V\mapsto M_\infty(V)$ 到本讲义的猜想，这一步大致是先由 $\mathrm{PGL}_2(\Z_p)$ 的表示走到 $\mathrm{PGL}_2(\Q_p)$ 的表示（把 $V$ 换成紧诱导 $c\text{-}\mathrm{Ind}_{\mathrm{PGL}_2(\Z_p)}^{\mathrm{PGL}_2(\Q_p)}V$），再由伽罗瓦形变环上的模走到 $(\varphi,\Gamma)$-模叠上的层。

实际上，在拼接的语境中，从 $\mathrm{PGL}_2(\Z_p)$ 走到 $\mathrm{PGL}_2(\Q_p)$ 是 [CEG+16] 的主要贡献之一；于是思想的推进很自然：[EGS15]（拼接作为 $\mathrm{PGL}_2(\Z_p)$-表示上的函子）→ [CEG+16]（拼接作为 $\mathrm{PGL}_2(\Q_p)$-表示上的函子）→ 下文的预期定理 7.2.1（拼接升级为从 $\mathrm{PGL}_2(\Q_p)$-表示到 $(\varphi,\Gamma)$-模叠上（复形的）层的函子）。

如上文已指出的，拼接方法可从模曲线推广到更一般的志村簇及相关同余商。不过此处考虑的 $\mathrm{PGL}_2(\Q_p)$ 情形特别好，因为此情形已有 $p$-进局部 Langlands 对应（实际上是 $\mathrm{GL}_2(\Q_p)$ 的）。$p$-进局部 Langlands 对应的想法由 Taylor--Wiles 拼接所启发，但对应本身是用相当不同的技术构造的。本节余下部分回顾其中一些历史。

\ysrc{185}
\subsection{$p$-进 Langlands：概览与历史}

继许多研究者（按大致的历史顺序：Fontaine--Langlands、Serre、Mazur、Harris、Vignéras 与 Schneider--Teitelbaum）初步的、大多未发表的研究之后，对可能的 $p$-进 Langlands 对应的研究——把群 $\mathrm{GL}_2(\Q_p)$ 的 $p$-进表示与局部伽罗瓦群 $\mathrm{Gal}_{\Q_p}$ 的二维 $p$-进表示联系起来——由 Christophe Breuil 发起 [Bre03a, Bre03b]。模性提升定理的证明过程中产生的各种考量在相当大程度上推动了他 [Wil95, TW95, CDT99, BCDT01]；特别地，直接推动他的是 Breuil--Mézard 猜想所引发的考量（关于这些动机的叙述见 [Bre08, §1.4]）。我们先讨论与这些定理相关的一些想法与方法，再讨论对应本身、它在 Breuil--Mézard 猜想研究中的应用，以及它经由局部--整体相容性与完备（上）同调的关系。

\subsubsection{Taylor--Wiles--Kisin 拼接}

如上文简要回顾的，[TW95] 引入的 Taylor--Wiles 拼接是一个从志村簇的（上）同调（视为 Hecke 模）构造局部形变环上的模的过程。

拼接方法最初是为证明模性提升定理而发展的，即证明某些 $p$-进伽罗瓦表示（在最初的应用中，指与椭圆曲线相伴者）对应于模形式。起初，拼接论证需要以若干关于模形式之间同余的深刻结果为输入，包括 mod $p$ 重数一定理、提高水平（level raising）定理、Serre 猜想权部分的若干情形，以及 Serre 模性猜想本身。此后局面反转：如今这类结果反倒可以从拼接论证推出。

拼接论证至今仍用来证明伽罗瓦表示是自守的，但如今也用来研究伽罗瓦表示以及自守表示之间同余的内部性质，我们感兴趣的正是后一应用。我们不打算完整叙述这些转变，而只讨论拼接构造的当代视角；简言之：对 mod $p$ 重数一定理的依赖由 Diamond [Dia97] 与 Fujiwara [Fuj06] 各自独立地消除，提高水平由 Taylor 的「Ihara 回避」（Ihara avoidance）论证 [Tay08] 处理，Serre 猜想的权部分由 T.G. [Gee11] 处理，Serre 猜想则由 Khare--Wintenberger [KW09] 与 Kisin [Kis09b] 证明。这些发展中的许多，以及下文所解释的大部分内容，都关键地依赖 Kisin 对 Taylor--Wiles 方法的改进 [Kis09c]。

我们沿用上文引入的记号，只作两处调整以免累赘：从今往后把 $R(\overline{\rho})$ 写作 $R_p$；并允许整体形变问题中固定的行列式为任意奇特征，而不只是 $\varepsilon^{-1}$，于是相应的局部形变问题中固定的行列式便允许是任意特征。允许局部行列式为任意特征的意义在于：这样下文便可考虑 $\mathrm{GL}_2(\F_p)$、$\mathrm{GL}_2(\Z_p)$ 或 $\mathrm{GL}_2(\Q_p)$ 的、具有固定但任意的中心特征的表示——这正是 $p$-进局部 Langlands 的通常设定。

\ysrc{186}
例如，在 $\mathrm{GL}_2$（而非 $\mathrm{PGL}_2$）的模曲线上作拼接，便得到上文所述正合函子 $V\mapsto M_\infty(V)$ 的一个版本，它把有限生成 $\OO$-模上、带适当指定中心特征的连续 $\mathrm{GL}_2(\Z_p)$-表示范畴，送到 $R_\infty$ 上的有限生成模范畴，其中 $R_\infty$ 是 $R_p$ 上的幂级数环。（此处及以下讨论中，我们总假设形变问题所固定的行列式与 $\mathrm{GL}_2$-表示的中心特征已作相容选取；由于固定行列式与中心特征这一问题主要是技术性的，我们尽量不赘述。）

有些 $\mathrm{GL}_2(\Z_p)$-表示上的 $M_\infty$ 取值格外令人感兴趣：其一为局部代数型（locally algebraic type）中的格，其二为 $\mathrm{GL}_2(\F_p)$ 在 $k$ 上的不可约表示。在第一类表示上计算拼接函子 $M_\infty$，把拼接与模性提升的研究联系起来；这正是拼接构造最初的动机。在第二类表示上计算拼接函子，则建立拼接与 Serre 猜想权部分的联系。

更确切地说，回忆所谓局部代数型 $\sigma$，指的是 $\mathrm{GL}_2(\Z_p)$ 的不可约 $\Q_p$-表示，由代数表示与一个经由 Henniart 的惯性局部 Langlands 对应 [BM02, App. A] 与惯性型 $\tau$ 相对应的有限维表示作张量而得；若 $\overline{\rho}$ 的一个提升具有与代数表示的最高权相对应的 Hodge--Tate 权、且惯性型对应于 $\tau$，则称它是 $\sigma$ 型的提升。设 $\sigma$ 是局部代数型，$\sigma^\circ$ 是 $\sigma$ 中的 $\mathrm{GL}_2(\Z_p)$-不变格。若 $R_p(\sigma)$ 记 $R_p$ 中参数化 $\sigma$ 型提升的商，则（作为局部--整体相容性的推论）$M_\infty(\sigma^\circ)$ 上的 $R_p$-作用经由 $R_p(\sigma)$ 分解；而关于 $\sigma$ 型提升的模性提升定理（即证明 $\overline{r}$ 的 $\sigma$ 型提升必为模的定理）用拼接的语言可以说成：$R_\infty$ 的相应商 $R_\infty(\sigma)$ 在 $M_\infty(\sigma)$ 上的作用应当是忠实的。\tn{底本此处作 $M_\infty(\sigma)$，但据上下文（$R_\infty(\sigma)$ 被定义为在 $M_\infty(\sigma^\circ)$ 上忠实作用的商，参本页下文与 §3.2.3）疑为 $M_\infty(\sigma^\circ)$ 之误，译文照原。}

至于 Serre 猜想的权部分，则可重述为：$M_\infty(\overline{\sigma})$ 非零，当且仅当 $\mathrm{GL}_2(\F_p)$ 的 Serre 权 $\overline{\sigma}$ 属于某个指定集合 $W(\overline{\rho})$（按定义，Serre 权是 $\mathrm{GL}_2(\F_p)$ 的不可约 $k$-表示）。

\subsubsection{一个观察}

若 $V$ 是有限秩自由 $\OO$-模上的 $\mathrm{GL}_2(\Z_p)$-表示（例如局部代数型中的一个格），且 $\overline{V}:=V/\varpi V$ 是相应的剩余表示，则（由正合性）$M_\infty(V)$ 是 $\OO$-无挠的，于是 $M_\infty(V)\neq0$ 当且仅当 $M_\infty(\overline{V})\neq0$，当且仅当对 $\overline{V}$ 的某个 Jordan--Hölder 因子（Jordan--Hölder factor）$\overline{\sigma}$ 有 $M_\infty(\overline{\sigma})\neq0$。因此，若 Serre 猜想的权部分已知，则 $M_\infty(V)$ 是否非零这一问题可以通过考察 $\overline{V}$ 的 Jordan--Hölder 因子集合来回答。这一简单的观察推动了该学科此后的大部分发展。

\subsubsection{Breuil--Mézard 猜想}

如前所述，要证明局部代数型 $\sigma$ 的伽罗瓦表示的模性提升定理，须证明 $R_\infty(\sigma)$ 在 $M_\infty(\sigma^\circ)$ 上忠实作用，其中 $\sigma^\circ$ 是 $\sigma$ 中某个（等价地，任一）不变格。由拼接构造可知，$M_\infty(\sigma^\circ)$ 在 $R_\infty(\sigma)$ 上是极大 Cohen--Macaulay 的；故若 $R_p(\sigma)$ 是整环，则 $R_\infty(\sigma)$（此时也是整环）必在 $M_\infty(\sigma^\circ)$ 上忠实作用。

\ysrc{187}
或许证明 $R_p(\sigma)$ 是整环最简单的办法是证明它形式光滑；[CDT99] 与 [BCDT01] 的论证正是为建立这样的光滑性结果（对适当选取的 $\sigma$）。Breuil 与 Mézard 在分析这两篇论文中为确认此类结果而引入的条件时，表述了他们那个同名的猜想 [BM02]。该猜想用 $\overline{\rho}$ 的、作为 $\sigma^\circ/\varpi$ 的成分出现的 Serre 权，来描述 $R_p(\sigma)/\varpi$ 的 Hilbert--Samuel 重数（Hilbert--Samuel multiplicity）。

Kisin 在理解 Breuil--Mézard 猜想上取得关键突破（例如见 [Kis10b]）：他强化了 §3.2.2 的观察，指出由 Serre 猜想的权部分、拼接函子的正合性以及拼接模的极大 Cohen--Macaulay 性可推出，Breuil--Mézard 猜想的公式恰好描述 $R/\varpi$ 的 Hilbert--Samuel 重数，其中 $R$ 记 $R_\infty(\sigma)$ 中在 $M_\infty(\sigma^\circ)$ 上忠实作用的那个商。于是，Breuil--Mézard 猜想与 Serre 猜想的权部分合起来，便本质等价于 $R_p(\sigma)$ 在 $M_\infty(\sigma^\circ)$ 上作用的忠实性，即等价于一个模性提升定理。（关于这一等价在更一般设定下更精确的表述，见 [EG14, §5]。）

把这一观察与运用 $p$-进局部 Langlands 对应的论证结合起来，Kisin 事实上证明了 Breuil--Mézard 猜想（并由此推出相应的模性提升定理）[Kis09a]。下文还会回顾这一论证的更多细节。目前只需指出：即便不知道 $R_p(\sigma)$ 是否在 $M_\infty(\sigma^\circ)$ 上忠实作用，Hilbert--Samuel 重数也有显然的不等式 $e_{R_p(\sigma)/\varpi}\ge e_{R/\varpi}$；因此 Kisin 的观察总能给出 Breuil--Mézard 猜想的一个弱形式，其中等号被替换为单向的不等式。

\subsubsection{局部解析的变体}

上文讨论了经由拼接函子研究 Serre 权的途径。解析设定中有相关的故事，见 [Bre15] 与 [BHS19]。哪些 Serre 权与剩余伽罗瓦表示 $\overline{\rho}$ 相关联，这一问题可理解为：哪些不可约 $\mathrm{GL}_2(\Z_p)$-表示（更一般地，$K=\mathrm{GL}_n(\OO_F)$-表示）嵌入（猜想中的）mod $p$ 表示 $\Pi(\overline{\rho})$。（注意此处我们采取的是关于 $p$-进与 mod $p$ 局部 Langlands 的更传统视角；这一视角与本讲义所强调的范畴猜想之间的联系，见注 7.8.6。）与此问题有某种类比，Breuil [Bre15, Conj. 6.1] 表述了一个猜想：局部解析主序列表示的哪些不可约成分嵌入 $p$-进伽罗瓦表示 $\rho$ 的 $p$-进局部 Langlands 对应 $\Pi(\rho)$。关于 Serre 权的问题由如下整体问题驱动：相伴伽罗瓦表示在模 $p$ 下有固定约化的模形式，其权为何；而 Breuil 的基座（socle）猜想则由如下整体问题驱动：具有指定相伴伽罗瓦表示的非经典超收敛 $p$-进模形式（其权）为何。

在 [BHS19] 中，Breuil 基座猜想的许多情形用类似策略（拼接函子，以及一个 Breuil--Mézard 猜想风格的计重公式）得到证明：我们不再把 $V\mapsto M_\infty(V)$ 作为某些局部代数 $K$-表示上的函子来研究，转而研究一个相关的函子，它作用于某些局部解析表示（或 $\mathcal{O}$ 范畴，以及从该范畴到局部解析表示的函子）。

\ysrc{188}
\subsubsection{$\mathrm{GL}_2(\Q_p)$ 的 $p$-进局部 Langlands 对应}

此时该函子产出 $\Spf R_\infty$ 的刚性解析一般纤维上的凝聚层，支撑在以 trianguline 伽罗瓦表示定义的闭子簇上（关于 trianguline 表示的讨论亦见下文 §3.2.6）。对这些 trianguline 形变空间，[BHS19, Conj. 4.3.4] 表述了 Breuil--Mézard 猜想的一个局部解析类比。粗略地说，该猜想预测：从特征 $\delta$ 抛物诱导所得的局部解析表示中，各不可约成分的重数，与参数对应于 $\delta$ 的 trianguline 形变空间中某些闭链的重数相匹配。事实上，通过把它与几何表示论中的计重公式联系起来，该猜想在许多情形可证 [BHS19, Theorem 4.3.8]。

$p$-进局部 Langlands 对应把 $\mathrm{GL}_2(\Q_p)$ 的容许连续酉 $p$-进 Banach 空间表示 $\Pi(\rho)$ 指派给连续二维表示 $\rho:\mathrm{Gal}_{\Q_p}\to\mathrm{GL}_2(L)$；Breuil 最先对潜在半稳定且 Hodge--Tate 权互异的表示 $\rho$ 提出此对应 [Bre03a, Bre03b, Bre04]。

更确切地说，Breuil 在此背景下提出，$\Pi(\rho)$ 应构造为某个局部代数表示 $\Pi(\rho)^{\mathrm{alg}}$ 的某种完备化；后者由与 $\mathrm{WD}(\rho)$ 相伴的 $\mathrm{GL}_2(\Q_p)$ 光滑表示（$\mathrm{WD}(\rho)$ 即潜在半稳定表示 $\rho$ 的 Dieudonné 模底下的 Weil--Deligne 表示）与一个编码 $\rho$ 的 Hodge--Tate 权的 $\mathrm{GL}_2(\Q_p)$ 有限维代数表示作张量而得。（表示 $\Pi(\rho)^{\mathrm{alg}}$ 含有局部代数型 $\sigma$，而 Breuil 的提议受 Breuil--Mézard 猜想驱动。）

此外，当 $\rho$ 是晶体阿贝尔型（crystabelline）或半稳定阿贝尔型（semistabelline）（即限制到 $\Q_p$ 的某个有限阿贝尔扩张后变为晶体或半稳定）时，他让提议完全显式化，即就 $\Pi(\rho)$ 应当是什么给出具体提议。也就是说，在不可约晶体阿贝尔情形，他提议 $\Pi(\rho)$ 应是 $\Pi(\rho)^{\mathrm{alg}}$ 的泛酉完备化（universal unitary completion，[Eme05] 意义下）；而在不可约真正的半稳定阿贝尔（即非晶体阿贝尔）情形，他提议 $\Pi(\rho)$ 应是 $\Pi(\rho)^{\mathrm{alg}}$ 的某种依赖于 $\rho$ 的 $\mathcal{L}$-不变量的完备化。（如下文所述，这一提议仍有不显式的一面：实际上并不显然如何以任何具体方式计算这些泛酉完备化。）

在不可约晶体阿贝尔情形提议以泛酉完备化为 $\Pi(\rho)$ 的候选，其动机在于：此情形下 $\rho$ 并未承载 $\Pi(\rho)^{\mathrm{alg}}$ 尚未承载的「额外」$p$-进 Hodge 理论数据（也就是说，此情形下没有 $\mathcal{L}$-不变量之类的东西；更确切地说，在 $D_{\mathrm{cris}}(\rho)$ 上，具有固定 Hodge--Tate 权、且使相应伽罗瓦表示不可约的弱容许滤过在同构意义下唯一）。表示 $\Pi(\rho)^{\mathrm{alg}}$ 决定（且被决定）于 $D_{\mathrm{cris}}(\rho)$ 与 Hodge--Tate 权；而泛酉完备化是不依赖任何额外选择的、唯一显而易见的完备化选项。

研究 Breuil 所猜想的对应时，最初的主要困难在于确立他所提议的完备化的任何非平凡性质；甚至并不显然这些完备化非零。这一局面有一个例外：当 $\rho$ 是由 Hecke 特征形式产生的整体二维表示 $r:\mathrm{Gal}_{\Q}\to\mathrm{GL}_2(L)$ 到 $\mathrm{Gal}_{\Q_p}$ 的限制时。此时完备上同调理论连同经典 Langlands 对应的局部--整体相容性，给出 $\Pi(\rho)^{\mathrm{alg}}$ 的一个非零酉完备化。在半稳定阿贝尔情形，Breuil 进一步用源于 $p$-进 $L$-函数理论的方法（并提醒读者：$\mathcal{L}$-不变量的概念正起源于同一理论）证明，该完备化经由他所提议的 $\Pi(\rho)^{\mathrm{alg}}$ 的完备化分解 [Bre10b]。

\ysrc{189}
Kato、Perrin-Riou 等人经由伽罗瓦上同调、显式互反律与 $p$-进 Hodge 理论处理 $p$-进 $L$-函数理论的途径启发了 Colmez；以他本人用 $(\varphi,\Gamma)$-模对这些想法所作的重释与进一步发展为基础，Colmez 取得了重大突破：在不可约且真正半稳定阿贝尔的情形，证明 Breuil 所提议的完备化 $\Pi(\rho)$ 非零、拓扑不可约且容许。（关于 Colmez 在这一方向上研究的高潮，见 [Col10a]。）

此后不久，Berger 与 Breuil 用类似技术证明了晶体阿贝尔情形的相应结果 [BB10]。他们结果的一个关键推论是：对晶体阿贝尔的 $\rho$，$\Pi(\rho)^{\mathrm{alg}}$ 的泛酉完备化 $\Pi(\rho)$ 是 $\Pi(\rho)^{\mathrm{alg}}$ 唯一的非零酉完备化。

由于 $(\varphi,\Gamma)$-模理论适用于任意表示 $\rho$，而不只是潜在半稳定者，Colmez 的结果使他提出：$p$-进局部 Langlands 对应应当对任意连续表示 $\rho:\mathrm{Gal}_{\Q_p}\to\mathrm{GL}_2(L)$ 存在。（当 $\rho$ 不是 Hodge--Tate 权互异的潜在半稳定表示时，人们不再有可用的表示 $\Pi(\rho)^{\mathrm{alg}}$，从而连猜想性地把 $\Pi(\rho)$ 构造成某种完备化都无从谈起。）

通过反向工程自己研究 Breuil 完备化的 $(\varphi,\Gamma)$-模途径，Colmez 发现了他的两个函子中的第一个：函子 $\Pi\mapsto V(\Pi)$，它把拓扑有限长的容许酉连续 Banach $\mathrm{GL}_2(\Q_p)$-表示（并满足其 mod $\varpi$ 约化为有限长这一条件，Paškūnas 证明该条件多余 [Paš13]）送到 $\mathrm{Gal}_{\Q_p}$ 的连续表示 [Col10b, §IV]。

鉴于自己对 Breuil--Mézard 猜想的研究，Kisin 随后提出了关键建议：Colmez 所提议的一般 $p$-进局部 Langlands 对应，应视为在局部伽罗瓦形变空间之上进行；形变论工具可在其构造中发挥作用。循此建议，Colmez 得以推广其 $(\varphi,\Gamma)$-模方法，对一般的 $\rho$ 构造 $\Pi(\rho)$ [Col10b]。两个函子经由自然同构 $V(\Pi(\rho))\cong\rho$ 相联系。此外他还证明：当 $\rho$ 是 Hodge--Tate 权互异的潜在半稳定表示时，表示 $\Pi(\rho)$ 确实是 $\Pi(\rho)^{\mathrm{alg}}$ 的完备化，与 Breuil 最初的提议相符。当 $\rho$ 是晶体阿贝尔或半稳定阿贝尔时，这一点可由上文提到的他本人以及 Berger--Breuil 的结果推出（事实上，$\Pi(\rho)$ 的一般构造正是以那些情形的构造为动机）。然而在一般情况下（即既非晶体阿贝尔、也非半稳定阿贝尔，更一般地，非 trianguline 的情形），$\mathrm{GL}_2(\Q_p)$ 整体在 $\Pi(\rho)$ 上的作用是用形变论方法（或者说，用刚性解析内插）以非常间接的方式构造的（只有 $\mathrm{GL}_2(\Q_p)$ 的上三角 Borel 子群的作用能直接用 $(\varphi,\Gamma)$-模看出），故 $\Pi(\rho)$ 中局部代数向量的性质并不显然。Colmez 关于它们包含（并事实上等同于）$\Pi(\rho)^{\mathrm{alg}}$ 的证明，归根结底靠与完备上同调的整体背景作比较，而这一比较以我们中一人（M.E.）关于局部--整体相容性的工作为中介。

\ysrc{190}
$\mathrm{GL}_2(\Q_p)$ 的 $p$-进局部 Langlands 对应理论由 Paškūnas 完成：他证明函子 $\Pi\mapsto V(\Pi)$ 在 $\mathrm{GL}_2(\Q_p)$ 的超奇异拓扑不可约容许酉连续 Banach 空间表示的同构类，与 $\mathrm{Gal}_{\Q_p}$ 的不可约连续二维表示之间，诱导一个双射 [Paš13]。事实上还建立了更多：本质上，$V$ 在 $\mathrm{GL}_2(\Q_p)$ 的拓扑有限长容许 Banach 空间表示范畴，与对象为 $\mathrm{Gal}_{\Q_p}$ 的二维连续表示的形变之范畴之间，诱导一个等价。此外，该等价可用 Morita 理论描述：$\mathrm{GL}_2(\Q_p)$-表示范畴中有某些内射对象，其自同态环可等同于伽罗瓦形变环，由它们充当中介。由于这些内射对象（等价地，其对偶——此时是适当范畴中的投射对象）在 Andrea Dotto、M.E. 与 T.G. 的工作（我们在 §7.2–7.4 中说明）中起重要作用，下文将进一步讨论它们。

\subsubsection{局部解析的 $p$-进 Langlands 对应}

如下文将进一步讨论的，$\mathrm{GL}_2(\Q_p)$ 的 $p$-进局部 Langlands 对应中出现的 Banach 空间表示，最终是经由取 mod $p^n$ 构造的极限（再倒置 $p$）而得到的；特别地，它们自然落入本讲义的「Banach」背景。\tn{底本此处作「further below below」，「below」重出，译文按一次「下文」处理。}

另一方面，研究 $\Pi(\rho)$ 这类 Banach 空间表示的一项重要工具，是研究其解析向量空间 $\Pi(\rho)^{\mathrm{an}}$（[ST03, §7] 意义下），它们是 $\mathrm{GL}_2(\Q_p)$ 的局部解析表示。这与考虑 mod $p^n$ 表示的方向相反：$\Pi(\rho)^{\mathrm{an}}$ 忘掉了 $\Pi(\rho)$ 上的 Banach 空间（从而 $p$-进整）结构。另一方面，Colmez 的构造表明，局部解析表示 $\Pi(\rho)^{\mathrm{an}}$ 可直接由与 $\rho$ 相伴的、Robba 环 $\mathcal{R}$ 上的刚性解析 $(\varphi,\Gamma)$-模 $D_{\mathrm{rig}}(\rho)$ 构造出来，故在此背景下研究 $p$-进局部 Langlands 对应仍有意义。这一观察推动 Colmez [Col16] 把 $\mathrm{GL}_2(\Q_p)$ 的 $p$-进局部 Langlands 对应推广到 Robba 环上（可能非 étale 的）秩二 $(\varphi,\Gamma)$-模与 $\mathrm{GL}_2(\Q_p)$ 的局部解析表示的框架。这自然落入本讲义的「解析」背景。

我们要指出：通常不容易从相伴的 $(\varphi,\Gamma)$-模 $D(\rho)$ 读出表示 $\rho$ 的 $p$-进 Hodge 理论性质，但从刚性解析变体 $D_{\mathrm{rig}}(\rho)$ 读出则容易得多。因此，Hodge 理论信息——例如 $D_{\mathrm{dR}}(\rho)$ 上的 Hodge 滤过（假设 $\rho$ 是 de Rham 的）——应当与局部解析表示 $\Pi(\rho)^{\mathrm{an}}$ 的性质相联系（而直接以 Banach 空间表示 $\Pi(\rho)$ 理解这一点应当不容易）。下面我们详述此点。

$p$-进 Langlands 纲领最初的动机之一，是增进对模曲线（或更一般的志村簇）上同调中局部--整体相容性的理解。我们将在下文讨论 $p$-进局部--整体……
